\clearpage\SourceLeaf{529}\renewcommand{\BookRunningTitle}{Remarks.}
\pdfbookmark[0]{Remarks.}{bemerkungen-second-impression}
\begin{center}\textbf{REMARKS.}\end{center}
\BookRule{.16\textwidth}
In this second volume of \textsc{Gauss}’s works I have included all his memoirs, essays and tables on higher arithmetic, insofar as they are not already covered by the seven sections of the \textit{Disqu. Arithm.}: the five memoirs published in the ‘\textit{Commentationes societatis regiae scientiarum Gottingensis}’ (in quarto), the reviews of these works of his own and of some works by others published in the ‘\textit{Göttingischen Gelehrten Anzeigen}’ (in octavo) (unsigned by \textsc{Gauss}, but confirmed as his through the records of the Göttingen University Library), and a selection from his manuscript remains.
\par In the second printing I have also added the tables ‘Circuli quadratura nova’, ‘On Calculating Logarithms’, ‘Quadratorum myrias prima’, ‘Indices of Primes in the Higher Number Domain’, ‘Auxiliary Table for Solving the Indeterminate Equation $A=fxx+gyy$ by the Method of Exclusion’, as well as ‘Sectio octava’, as far as it was written down, and finally two letters from \textsc{Gauss} to \textsc{Dirichlet}, as essential parts of the history of higher arithmetic.
\par To make the subjects in so extensive a volume easier to survey, the theorems are distinguished typographically in the same way. To make both the earlier editions and the present one easier to use, references to the Disq. Arithm. give article numbers in place of page numbers, and references to memoirs give their own titles in place of their places of publication. The note originally attached to Article 2 of the memoir ‘\textit{Theorematis arithmetici demonstratio nova}’,
